\VerdierPhysicalPageAnchor{0024}%
deduces for the cohomology objects. The mechanism of spectral sequences most
often conceals the simplicity of the phenomena. The elementary example of
the Künneth relations is significant in this respect.

\VerdierSemanticAnchor{intro-5}
\noindent\textbf{5.}\quad
Let $X$ and $Y$ be two compact spaces, let $A$ be a commutative ring with
identity (for simplicity), and let $F$ and $G$ be two sheaves of $A$-modules
on $X$ and $Y$, respectively. When $A$ is a field, it is known that the cross
product [23] defines an isomorphism of graded vector spaces:

\VerdierSemanticAnchor{intro-5-kunneth-field}
\begin{equation}
\mathrm H^*(X,F)\otimes_A\mathrm H^*(Y,G)
\xrightarrow{\sim}
\mathrm H^*(X\times Y,F\otimes_A G).
\tag{$\text{Kü}_{\mathrm{coh}}^{1}$}
\end{equation}

When $A=\mathbb Z$ and either the sheaf $F$ or the sheaf $G$ is torsion-free,
one has split exact sequences:

\VerdierSemanticAnchor{intro-5-kunneth-torsion-free}
\begin{equation}
\begin{aligned}
0 &\longrightarrow
\bigoplus_{p+q=n}\mathrm H^p(X,F)\otimes_A\mathrm H^q(Y,G)
\longrightarrow \mathrm H^n(X\times Y,F\otimes_A G) \longrightarrow \\
&\longrightarrow
\bigoplus_{p+q=n+1}
\operatorname{Tor}^{\mathbb Z}\bigl(\mathrm H^p(X,F),\mathrm H^q(Y,G)\bigr)
\longrightarrow 0.
\end{aligned}
\tag{$\text{Kü}_{\mathrm{coh}}^{2}$}
\end{equation}

When $A=\mathbb Z$ and no assumption is made on $F$ and $G$, one has two
families of split exact sequences:

\VerdierSemanticAnchor{intro-5-kunneth-general-split}
\begin{equation}
\begin{aligned}
0 &\longrightarrow
\bigoplus_{p+q=n}\mathrm H^p(X,F)\otimes_A\mathrm H^q(Y,G)
\longrightarrow \mathbb H^n \longrightarrow \\
&\longrightarrow
\bigoplus_{p+q=n+1}
\operatorname{Tor}^{\mathbb Z}\bigl(\mathrm H^p(X,F),\mathrm H^q(Y,G)\bigr)
\longrightarrow 0, \\
0 &\longrightarrow
\mathrm H^{n+1}\bigl(X\times Y,\mathcal T\!or^A(F,G)\bigr)
\longrightarrow \mathbb H^n
\longrightarrow \mathrm H^n(X\times Y,F\otimes_A G)
\longrightarrow 0.
\end{aligned}
\tag{$\text{Kü}_{\mathrm{coh}}^{3}$}
\end{equation}

The initial formulation $\text{(Kü)}_{\mathrm{coh}}^{1}$ thus
becomes more complicated, but remains accessible to explicit calculation.
When the ring $A$ is more complicated, for example when
$A=\mathbb Z/\ell^r\mathbb Z$, with $\ell$ prime and $r>1$, a ring of infinite
global cohomological dimension, traditionally one has only two spectral
sequences with the same abutment, whose initial terms are:

\VerdierSemanticAnchor{intro-5-kunneth-spectral-sequences}
\begin{equation}
\begin{aligned}
{}'K_2^{p,q}
  &=\bigoplus_{r+s=q}
    \operatorname{Tor}_{-p}^{A}
    \bigl(\mathrm H^r(X,F),\mathrm H^s(Y,G)\bigr), \\
{}''K_2^{p,q}
  &=\mathrm H^p\bigl(X\times Y,\mathcal T\!or_{-q}^{A}(F,G)\bigr).
\end{aligned}
\tag{$\text{Kü}_{\mathrm{coh}}^{4}$}
\end{equation}

However, these spectral sequences and the isomorphism between their
abutments are only the consequence and imperfect translation of the
following relation in the derived category of the category of $A$-modules:

\VerdierSemanticAnchor{intro-5-kunneth-derived}
\begin{equation}
\mathrm{RH}(X,F)\overset{\mathrm L}{\otimes}_A\mathrm{RH}(Y,G)
\xrightarrow{\sim}
\mathrm{RH}\bigl(X\times Y,F\overset{\mathrm L}{\otimes}_A G\bigr),
\tag{$\text{Kü}_{\mathrm{Der}}$}
\end{equation}
